\section{Problema}

O Modelo de Atores é uma teoria matemática estabelecida por Carl Hewitt em 1973, baseada em trabalhos anteriores de Peter Bishop e Richard Steiger. Seus objetivos eram abordar os novos desafios impostos pelos sistemas massivamente concorrentes e paralelos que emergiram com a computação em nuvem e arquiteturas many-core.

Sua principal hipótese é que toda computação fisicamente possível pode ser diretamente implementada com Atores. O Modelo de Atores representa uma reformulação completa da computação concorrente, passando de máquinas de estado globais para sistemas distribuídos, assíncronos e baseados em passagem de mensagens que melhor refletem a realidade física dos sistemas de computação modernos.

Tradicionalmente, o Modelo de Atores tem sido amplamente usado na arquitetura de sistemas distribuídos, com implementações bem-sucedidas como Akka (Java), Orleans (C\#) e, mais notavelmente, Elixir/Erlang. No entanto, sua aplicação em aplicações centralizadas (não distribuídas) permanece relativamente inexplorada.
